\chapter{Mathematics}

OmniLaTeX combines \texttt{amsmath} and \texttt{unicode-math} to provide
a modern math typesetting system with OpenType math fonts, auto-scaling
delimiters, and convenience macros for statistics, physics, and
engineering notation.

% =========================================================================
\section{Equation Environments}
% =========================================================================

All standard \texttt{amsmath} environments are available.  The most
commonly used are summarised below with examples.

\textbf{Single equation:}

\begin{verbatim}
\begin{equation}
  E = mc^2
  \label{eq:einstein}
\end{equation}
\end{verbatim}

\textbf{Aligned multi-line:}

\begin{verbatim}
\begin{align}
  \nabla \cdot  \mathbf{E} &= \frac{\rho}{\varepsilon_0}  \label{eq:gauss} \\
  \nabla \times \mathbf{B} &= \mu_0 \mathbf{J}
    + \mu_0 \varepsilon_0 \frac{\partial \mathbf{E}}{\partial t} \label{eq:ampere}
\end{align}
\end{verbatim}

\textbf{Gathered (no alignment):}

\begin{verbatim}
\begin{gather}
  a^2 + b^2 = c^2 \\
  e^{i\pi} + 1 = 0
\end{gather}
\end{verbatim}

\textbf{Multi-line with single number:}

\begin{verbatim}
\begin{multline}
  \int_{-\infty}^{\infty} e^{-x^2} \, dx
    = \sqrt{\pi} \\
  = \sum_{n=0}^{\infty} \frac{(-1)^n}{n!}
    \int_{-\infty}^{\infty} x^{2n} e^{-x^2} \, dx
\end{multline}
\end{verbatim}

\textbf{Split inside equation:}

\begin{verbatim}
\begin{equation}
  \begin{split}
    f(x) &= a_0 + a_1 x + a_2 x^2 \\
         &= \sum_{k=0}^{2} a_k x^k
  \end{split}
\end{equation}
\end{verbatim}

% =========================================================================
\section{Auto-Scaling Delimiters}
% =========================================================================

Manually writing \cs{left(} and \cs{right)} for every expression is
tedious.  The \texttt{omnilatex-math} module provides three convenience
macros that automatically size delimiters based on their content:

\begin{table}[htbp]
  \centering
  \caption{Auto-scaling delimiter commands.}
  \label{tab:delimiters}
  \begin{tabular}{@{}lll@{}}
    \toprule
    Command                              & Delimiters & Example input         \\
    \midrule
    \texttt{\textbackslash parens\{x+y\}}  & $(\;)$ & \texttt{\textbackslash parens\{\textbackslash frac\{a\}\{b\}\}} \\
    \texttt{\textbackslash brackets\{x\}}  & $[\;]$ & \texttt{\textbackslash brackets\{\textbackslash Sigma\}} \\
    \texttt{\textbackslash braces\{x\}}    & $\{\;\}$ & \texttt{\textbackslash braces\{\textbackslash Omega\}} \\
    \bottomrule
  \end{tabular}
\end{table}

These commands expand to \cs{left} / \cs{right} pairs internally, so they
handle nested fractions, summations, and matrices without manual height
specification:

\begin{verbatim}
\parens{\frac{1}{1 + \frac{1}{n}}}
% → ( 1 / (1 + 1/n) )  with correctly sized parentheses
\end{verbatim}

% =========================================================================
\section{Equation Punctuation}
% =========================================================================

Placing punctuation after display equations is a common source of spacing
errors.  OmniLaTeX defines three macros that handle proper spacing
between the equation and the following punctuation:

\begin{table}[htbp]
  \centering
  \caption{Equation punctuation commands.}
  \label{tab:eq-punct}
  \begin{tabular}{@{}llp{0.55\textwidth}@{}}
    \toprule
    Command          & Produces & Usage \\
    \midrule
    \texttt{\textbackslash eqend}   & \eqend   & Period after a concluding equation \\
    \texttt{\textbackslash eqcomma} & \eqcomma & Comma before a clause continues \\
    \texttt{\textbackslash coloneq} & \coloneq & Colon for a definition equation \\
    \bottomrule
  \end{tabular}
\end{table}

Example:

\begin{verbatim}
\begin{equation}
  \mathcal{L}\{f(t)\} = \int_0^\infty e^{-st} f(t) \, dt \eqend
\end{equation}
The Laplace transform converts a time-domain signal\ldots
\end{verbatim}

% =========================================================================
\section{Statistical Macros}
% =========================================================================

The \texttt{omnilatex-math} module provides shorthand macros for common
statistical notation:

\begin{table}[htbp]
  \centering
  \caption{Statistical convenience macros.}
  \label{tab:stat-macros}
  \begin{tabular}{@{}lll@{}}
    \toprule
    Command                                 & Renders as & Meaning \\
    \midrule
    \texttt{\textbackslash mean\{x\}}      & $\mean{x}$      & Arithmetic mean (bar) \\
    \texttt{\textbackslash logmean\{x\}}   & $\logmean{x}$   & Logarithmic mean \\
    \texttt{\textbackslash abs\{x\}}       & $\abs{x}$       & Absolute value \\
    \texttt{\textbackslash norm\{x\}}      & $\norm{x}$      & Norm (double bars) \\
    \texttt{\textbackslash inner\{a,b\}}   & $\inner{a}{b}$  & Inner product (angle brackets) \\
    \texttt{\textbackslash variance\{x\}}  & $\variance{x}$  & Variance ($\sigma^2$) \\
    \bottomrule
  \end{tabular}
\end{table}

These macros use \cs{DeclarePairedDelimiter} internally, so they accept
optional size arguments: \texttt{\textbackslash mean*[\textbackslash Big]\{x\}} produces
a larger bar.

% =========================================================================
\section{Number Formatting}
% =========================================================================

\textbf{Circled numbers} are useful for step-by-step procedures and
enumerated equations:

\begin{verbatim}
\circlednum{1}  \circlednum{2}  \circlednum{3}
% → ①  ②  ③
\end{verbatim}

\textbf{Symbol placeholders} for figures and tables under development:

\begin{verbatim}
\symbolplaceholder{figure}
% → [FIGURE]
\end{verbatim}

\textbf{Custom operators} can be declared with \cs{DeclareMathOperator}:

\begin{verbatim}
\DeclareMathOperator{\Tr}{Tr}
\DeclareMathOperator*{\argmin}{arg\,min}

$\Tr(A) = \sum_i \lambda_i$
$\argmin_{x} f(x)$
\end{verbatim}

% =========================================================================
\section{Physical Constants}
% =========================================================================

OmniLaTeX defines shortcuts for frequently used constants and symbols in
engineering and physics:

\begin{table}[htbp]
  \centering
  \caption{Physical constant macros.}
  \label{tab:constants}
  \begin{tabular}{@{}lll@{}}
    \toprule
    Command          & Renders as & Description \\
    \midrule
    \texttt{\textbackslash grad}  & $\grad$  & Gradient operator ($\nabla$) \\
    \texttt{\textbackslash const} & $\const$ & Denotes a constant quantity \\
    \bottomrule
  \end{tabular}
\end{table}

\textbf{Quoted text inside math mode} is provided by \cs{qq}\texttt{\{text\}}:

\begin{verbatim}
$T_{\qq{eff}} = 5778\,\mathrm{K}$
% → T_eff with "eff" in upright text
\end{verbatim}

This avoids the common pattern of switching to \cs{mathrm} or
\cs{textnormal} inline and keeps the source readable.

% =========================================================================
\section{Display Equations}
% =========================================================================

The \texttt{empheq} package allows annotated display equations with
side-by-side text or boxed emphasis.  OmniLaTeX preconfigures
\texttt{empheq} to work with the Libertinus Math font metrics:

\textbf{Boxed equation:}

\begin{verbatim}
\begin{empheq}[box=\fbox]{equation}
  i\hbar \frac{\partial}{\partial t} \Psi
    = \hat{H} \Psi
\end{empheq}
\end{verbatim}

\textbf{Equation with right-side annotation:}

\begin{verbatim}
\begin{empheq}[right=\eqref{eq:schrodinger}\;\text{(Schr\"odinger)}]{align}
  i\hbar \frac{\partial}{\partial t} \Psi
    &= \hat{H} \Psi \label{eq:schrodinger}
\end{empheq}
\end{verbatim}

\textbf{Aligned equations with shared numbering:}

\begin{verbatim}
\begin{align}
  \frac{\partial u}{\partial t}
    &= \alpha \nabla^2 u \label{eq:heat} \\
  \frac{\partial^2 u}{\partial t^2}
    &= c^2 \nabla^2 u   \label{eq:wave}
\end{align}
\end{verbatim}

Cross-referencing equations works exactly like \LaTeX\ figures and
tables: \texttt{Equation~\textbackslash ref\{eq:heat\}} produces
\emph{Equation~1}.
